\uuid{yzbP}
\exo7id{822}
\titre{exo7 822}
\auteur{bodin}
\organisation{exo7}
\datecreate{1998-09-01}
\isIndication{false}
\isCorrection{false}
\chapitre{Calcul d'intégrales}
\sousChapitre{Polynôme en sin, cos ou en sh, ch}
\cpge{mpsi-mp2i.primitives-equations-differentielles.a, pcsi.techniques-calcul-primitives.a, ptsi.techniques-calcul-primitives.a, mpsi-mp2i.calcul-algebrique-trigonometrie.d, pcsi.calcul-algebrique-trigonometrie.d, ptsi.calcul-algebrique-trigonometrie.d, bcpst1.derivees-primitives-integrales.c, bcpst1.derivees-primitives-integrales.d}
\module{Analyse}
\niveau{L1}
\difficulte{}

\contenu{
\question{
Calculer les primitives suivantes :
$$
\int (\cos x \cos 2x + \sin x \sin 3x)dx \quad ; \quad
\int \cos x \sin^4x dx \quad ; \quad
\int \cos^6x dx \quad ;
$$

$$
\int \sin^3x \cos x dx \quad ; \quad
\int \sin^4x dx \quad ; \quad
\int \sin^3x \cos^2x dx \quad ;
$$

$$
\int \ch^2 x \sh^2 x dx \quad ; \quad
\int \sh x \ch^3 x dx \quad ; \quad
\int \ch x \sh^3x dx.
$$
}
}
